%% ==============================
\chapter{Stand der Wissenschaft und Technik}
\label{cha:state_of_the_art}
%% ==============================

Ziel dieses Kapitels ist es, den aktuellen Forschungsstand zu dezentralen DLT-basierten Architekturen für Gesundheitsdatenmanagement darzustellen. Der Fokus liegt auf blockchainbasierten sowie hashgraphbasierten Ansätzen und deren Einsatz im Gesundheitswesen. Dazu werden bestehende Ansätze aus der Literatur analysiert und bezüglich ihrer Anforderungen, Architekturbausteine, Herausforderungen und Vergleichskriterien eingeordnet. Außerdem wird auch die Literatur zu konkreten technischen Frameworks der Distributed-Ledger-Technologien untersucht.

Das Kapitel gliedert sich wie folgt: Abschnitt~\ref{sec:literaturrecherche} beschreibt die Methodik der durchgeführten systematischen Literaturrecherche, einschließlich Suchstring und Auswahlprozess. Abschnitt~\ref{sec:uebersicht_literatur} stellt die identifizierte Literatur inhaltlich vor, gegliedert in konkrete DLT-Architekturen im Gesundheitswesen, vergleichende und methodische Arbeiten sowie ergänzende Literatur außerhalb des Gesundheitskontexts. Abschnitt~\ref{sec:literatur_architekturvergleich} ordnet die beschriebenen Architekturen taxonomisch ein, und Abschnitt~\ref{sec:literatur_schlussfolgerungen} leitet daraus die Forschungslücke sowie die Vergleichskriterien dieser Arbeit ab. Der Abschnitt~\ref{sec:sota_referenzarchitektur} beschreibt abschließend die Referenzarchitektur im Detail, deren Ledger-Schicht den Gegenstand des in dieser Arbeit durchgeführten Vergleichs bildet.


\section{Literaturrecherche}
\label{sec:literaturrecherche}
Um dieses Ziel zu erreichen, wurde eine strukturierte Literaturrecherche durchgeführt. Zuerst wurden die Schlüsselwörter (Keywords) gesucht und in einen logischen Zusammenhang gebracht; daraus wurde eine Wortliste aus Synonymen, Oberbegriffen, Unterbegriffen und verwandten Begriffen zusammengestellt, um ein möglichst großes Spektrum abzudecken. Danach wurde durch AND- und OR-Verknüpfungen aus mehreren Gruppen von Schlüsselwörtern ein Suchstring (Search String) gebildet. Mithilfe des gebildeten Suchstrings wurden die wissenschaftlichen Datenbanken Scopus, IEEE und PubMed durchsucht. Des Weiteren wurden die Ergebnisse der Suche analysiert und aussortiert, bis eine Liste der relevanten Literatur fertiggestellt war.
„Blockchain-based" wurde statt „blockchain" verwendet, um irrelevante Papers zu vermeiden. Die verwendeten Schlüsselwörter gliedern sich in fünf Gruppen: Für den Gesundheitskontext wurden die Begriffe EHR, electronic health record, medical data und patient records verwendet; für die zugrunde liegende DLT die Begriffe DLT, distributed ledger, blockchain-based und hashgraph; für die Architekturebene die Begriffe architecture, data management, data sharing und framework; als Evaluationskriterien performance, scalability, security, throughput, fault tolerance, latency und evaluation; und als Ausschlusskriterien supply chain, cryptocurrency, bitcoin, finance, insurance und IoT.

Für die Suche in den wissenschaftlichen Datenbanken wird ein Suchstring verwendet, der aus mehreren Blöcken besteht. Die Schlüsselwörter innerhalb der Blöcke werden durch OR verbunden, die Blöcke selbst durch AND. Zusätzlich wurden mittels NOT-Verknüpfung für diese Arbeit nicht relevante Bereiche wie Supply-Chain-Management, Kryptowährungen, Finanzwesen, Versicherungen und IoT-Anwendungen ausgeschlossen. Als Suchfeld wurde TITLE-ABS-KEY (bzw. die entsprechende Filter-Option der jeweiligen Datenbank) ausgewählt, da die relevanten Begriffe zu den Artikeln entweder im Titel, im Abstract oder in den Schlüsselwörtern zu finden sein sollten.

\begingroup
\shorthandoff{"}
\begin{quote}
\raggedright
\ttfamily
("electronic health record" OR EHR OR "medical data" OR "patient records")
AND ("distributed ledger" OR DLT OR "blockchain-based" OR hashgraph) 
AND (architecture OR framework OR "data sharing" OR "data management") 
AND (performance OR scalability OR security OR throughput OR latency OR "fault tolerance" OR evaluation)
AND NOT ("supply chain" OR cryptocurrency OR bitcoin OR finance OR insurance OR IoT)
\end{quote}
\endgroup

Durch die Suche in den wissenschaftlichen Datenbanken mithilfe des abgeleiteten Suchstrings wurden bei IEEE 276 Treffer, bei Scopus 747 Treffer und bei PubMed 63 Treffer gefunden. Nach dem Entfernen der Duplikate blieben 814 Papers übrig. Diese wurden anschließend anhand des Titels auf 105 Kandidaten für relevante Paper reduziert. Danach wurden die Abstracts gelesen und nur noch 43 Papers zum Lesen des Volltextes ausgewählt. Nach dem Lesen des Volltextes wurden 19 relevante Paper ausgewählt: 14 Arbeiten, die eine konkrete Architektur vorschlagen und evaluieren (Abschnitt~\ref{sec:literatur_architectures}), sowie fünf vergleichende und methodische Arbeiten (Abschnitt~\ref{sec:literatur_vergleichsmethodik}).

Da der oben definierte Suchstring auf den Gesundheitskontext beschränkt war, wurde zusätzlich gezielt außerhalb dieses Suchstrings nach Literatur recherchiert, die hashgraphbasierte SSI/DID-Mechanismen unabhängig vom Gesundheitswesen untersucht. Der Grund dafür war die im Abschnitt~\ref{sec:literatur_schlussfolgerungen} begründete Beobachtung, dass die gesundheitsspezifische Literatur keine Architektur mit einem hashgraphbasierten Ledger und einem SSI/DID-Identitätsmodell kombiniert; die zusätzliche Recherche prüft, ob eine solche Kombination zumindest außerhalb des Gesundheitskontexts existiert. Für die ergänzende Recherche wurde nach demselben Prinzip ein zweiter Suchstring gebildet, der die SSI/DID-Terminologie mit hashgraphspezifischen Begriffen kombiniert. Als Suchfeld diente wie zuvor TITLE-ABS-KEY.

\begingroup
\shorthandoff{"}
\begin{quote}
\raggedright
\ttfamily
("self-sovereign identity" OR SSI OR "identity management" OR "decentralized identifier" OR DID OR "verifiable credential")
AND (hashgraph OR hedera OR HCS)
\end{quote}
\endgroup

Die Suche ergab bei IEEE 103, bei Scopus 12 und bei PubMed 38 Treffer. Nach dem Entfernen der Duplikate verblieben 145 Papers, von denen anhand des Titels 11 Kandidaten für relevante Literatur ausgewählt wurden. Nach dem Lesen der Abstracts und Volltexte wurden schließlich zwei Arbeiten als relevant eingestuft, die in Abschnitt~\ref{sec:literatur_ergaenzend} als ergänzende Literatur beschrieben werden: Satybaldy et al. \cite{satybaldy2026} sowie Cádiz et al. \cite{cadiz2025}.


\section{Übersicht der relevanten Literatur}
\label{sec:uebersicht_literatur}
Dieses Kapitel fasst die Arbeiten zusammen, die im Rahmen der systematischen Literaturrecherche identifiziert und als relevant eingestuft wurden, und ordnet sie in den fachlichen Kontext dieser Arbeit ein. Während die grundlegenden technologischen Konzepte (u.\,a. Blockchain, Hashgraph und SSI/DID) bereits im Grundlagenkapitel~\ref{cha:grundlagen} vorgestellt wurden, liegt der Fokus der folgenden Abschnitte ausschließlich auf verwandten Arbeiten, die einen direkten inhaltlichen Bezug zur Zielsetzung dieser Arbeit aufweisen: der Konzeption, Implementierung und Evaluation dezentraler DLT-Architekturen für das Gesundheitsdatenmanagement. Die als Referenzarchitektur dienende Arbeit von Erler et al. \cite{erler2023} wird dabei zunächst wie alle übrigen Arbeiten in den Literaturblock eingeordnet und anschließend im Abschnitt~\ref{sec:sota_referenzarchitektur} im Detail beschrieben.

Die identifizierte Literatur kann anhand ihres jeweiligen Beitrags in drei thematisch und methodisch unterscheidbare Blöcke gegliedert werden. Der erste Block (Abschnitt~\ref{sec:literatur_architectures}) umfasst Arbeiten, die eine konkrete DLT-basierte Architektur oder ein Framework für das Gesundheitswesen vorschlagen und evaluieren. Diese werden nach dem verwendeten Identitätsmodell gegliedert, das der Unterscheidung von Erler et al. \cite{erler2024} folgt: SSI-basierte, DTI-basierte sowie schlüssel- bzw. zertifikatsbasierte Architekturen und Architekturen ohne ausgewiesenes Identitätsmodell. Die zugrunde liegende DLT-Familie wird innerhalb jeder Gruppe benannt, bildet jedoch bewusst nicht die Gliederungsebene, da sie erst in der Überlagerung mit dem Identitätsmodell aussagekräftig wird (Abschnitt~\ref{sec:literatur_schlussfolgerungen}). Jede Arbeit wird dabei so beschrieben, dass über alle Arbeiten hinweg eine konsistente Abstraktionsebene sichergestellt ist. Der zweite Block (Abschnitt~\ref{sec:literatur_vergleichsmethodik}) umfasst vergleichende und methodische Literatur in Form von Reviews und Benchmark-Studien, die selbst keine eigene Architektur vorschlagen, jedoch bestehende Systeme systematisch gegenüberstellen und dabei Bewertungskriterien, Metriken und offene Herausforderungen hervorheben. Der dritte Block (Abschnitt~\ref{sec:literatur_ergaenzend}) umfasst ergänzende Arbeiten außerhalb des Gesundheitskontexts, die hashgraphbasierte Ledger mit dezentralen Identitätsmechanismen verbinden.

Aufbauend auf dieser Übersicht werden die im Abschnitt~\ref{sec:literatur_architectures} beschriebenen Architekturansätze im Abschnitt~\ref{sec:literatur_architekturvergleich} taxonomisch eingeordnet. Aus dieser Taxonomie werden in Abschnitt~\ref{sec:literatur_schlussfolgerungen} die zentrale Forschungslücke sowie die für diese Arbeit abgeleiteten Vergleichskriterien entlang der Dimensionen Performanz, Skalierbarkeit und Sicherheit abgeleitet.

\subsection{Konkrete DLT-Architekturen im Gesundheitswesen}
\label{sec:literatur_architectures}
Der erste Literaturblock umfasst Arbeiten, die konkrete DLT-basierte Architekturen für das Gesundheitsdatenmanagement vorschlagen und bewerten. Gliederungskriterium ist das Identitätsmodell, da es bestimmt, wie anschlussfähig eine Architektur an die hier betrachtete Referenzarchitektur ist. Die vollständige Einordnung aller Arbeiten entlang von DLT-Art, Netzwerktyp, Framework, Konsens-Familie und Identitätsmodell leistet die Taxonomie in Abschnitt~\ref{sec:literatur_architekturvergleich}; dieser Abschnitt beschränkt sich darauf, je Gruppe das verbindende Architekturmuster und den jeweiligen Evaluationsansatz zu benennen.

\paragraph{Architekturen mit SSI-basiertem Identitätsmodell}
\label{par:literatur_ssi}

Erler et al.~\cite{erler2023} schlagen eine Architektur für den dezentralen Austausch medizinischer Dokumente zwischen Patient:innen und Einrichtungen vor, die auf einer Threat-Modeling-Analyse beruht. Sie kombiniert eine public permissioned SSI-Blockchain auf Basis von Hyperledger Indy (RBFT-Konsens über Indy Plenum) als Vertrauensanker mit einer DIDComm-basierten Kommunikationsschicht aus Hyperledger-Aries-Agenten. Der Ledger verwaltet ausschließlich pseudonyme Identitätsdaten -- DIDs, Schemas, Credential Definitions --, während medizinische Dokumente und Metadaten off-chain über Connector-Komponenten, Mediator-Agenten und eine Metadaten-Datenbank ausgetauscht werden~\cite{erler2023}. Eine Quantifizierung erfolgt nicht; validiert wird qualitativ gegen die aus dem Threat-Modeling abgeleiteten Sicherheitsanforderungen~\cite{erler2023}. Diese Arbeit dient als Referenzarchitektur und wird in Abschnitt~\ref{sec:sota_referenzarchitektur} im Detail dargestellt.

Drei weitere Arbeiten teilen dieses Grundmuster und unterscheiden sich in der Ausgestaltung. Manoj et al.~\cite{manoj2022} liefern die detaillierteste technische Beschreibung der Indy/Aries-Mechanismen -- paarweise eindeutige DIDs, Verifiable Credentials und die Indy-Rollen Steward, Trust Anchor und Endorser -- und messen Latenzen des Ausstellungs- und Verifikationsablaufs. Torongo und Toorani~\cite{torongo2023} bauen darauf auf und ergänzen ein dezentrales Schlüsselmanagement sowie einen expliziten Widerrufsmechanismus über kryptografische Akkumulatoren; ihre Evaluation verbindet Lasttests der Agenten-Endpunkte mit einer Sicherheitsanalyse nach dem STRIDE-Modell. Dieser Widerrufsmechanismus ist für die vorliegende Arbeit unmittelbar relevant, da seine Abbildung auf den Hedera Consensus Service eine der offenen Entwurfsfragen darstellt. Chainarong et al.~\cite{chainarong2025} verlassen die Indy-Familie und setzen DIDs und Verifiable Credentials auf eine regional gesharde Hyperledger-Fabric-Blockchain, kombiniert mit attributbasierter Verschlüsselung und Zero-Knowledge-Proofs; die Arbeit zeigt, wie SSI-Prinzipien mit Skalierbarkeitsmechanismen verknüpft werden können.

\paragraph{Architekturen mit DTI-basiertem Identitätsmodell}
\label{par:literatur_dti}

Bittins et al.~\cite{bittins2021} verbinden SSI/DID-Mechanismen mit etablierten Standards des Gesundheitswesens: Der Dokumentenaustausch läuft off-chain über die IHE-Profile XDS/XCA, während eine Hyperledger-Fabric-Schicht auf der europäischen EBSI-Infrastruktur Provenance-Nachweise und Credential-Prüfungen verankert. Eine quantitative Evaluation erfolgt nicht; validiert wird qualitativ gegen regulatorische und Interoperabilitätsanforderungen~\cite{bittins2021}. Relevant ist die Einordnung, dass SSI-Mechanismen als Ergänzung und nicht als Ersatz bestehender Austauschinfrastrukturen zu entwickeln sind.

\paragraph{Architekturen mit schlüssel- bzw. zertifikatsbasiertem Identitätsmodell}
\label{par:literatur_schluessel}

Fünf Arbeiten verzichten auf DIDs und stützen die Identität auf Zertifikate oder Schlüsselmaterial. Vier von ihnen verwenden permissioned Blockchains mit zertifikatsbasierter Mitgliederverwaltung und lagern die Akten verschlüsselt aus: Tawfik et al.~\cite{tawfik2025} trennen Speicher-, Kontroll- und Zugriffsebene und evaluieren systematisch mit einem Benchmark-Werkzeug entlang von Durchsatz, Latenz und Verhalten bei simulierten Angriffen -- ein methodischer Maßstab, an dem sich das Evaluationskonzept dieser Arbeit orientiert. Joshi und Mahajan~\cite{joshimahajan2025} stellen dazu eine direkte Konsens-Alternative dar und begründen ihre Wahl mit Skalierbarkeits- und Fehlertoleranzanforderungen. Mamoru et al.~\cite{mamoru2026} etablieren den Widerruf als eigenständiges, messbares Kriterium und erheben neben der Latenz von Einwilligungs- und Widerrufsoperationen auch die Durchsetzungsquote widerrufener Zugriffe; diese Größe ist unmittelbar auf die Frage übertragbar, wie sich die Revocation Registry der Referenzarchitektur auf HCS abbilden lässt. Farin et al.~\cite{farin2025} ergänzen einen kriterienbasierten Systemvergleich entlang von Konsens, Netzwerktyp, Speichermodell, Zugriffskontrolle, Transaktionskosten und Skalierbarkeit, der als Vorbild für die Vergleichsdimensionen dieser Arbeit dient. Minango et al.~\cite{minango2025} schließlich verwenden als einzige Arbeit dieser Gruppe keine Blockchain, sondern Corda, bei dem Transaktionsdaten nur zwischen den beteiligten Knoten geteilt werden und ein Notary-Dienst die Eindeutigkeit von Zuständen sichert; ihre Evaluation erhebt neben Durchsatz und Latenz auch CPU- und Speicherauslastung und erweitert das Spektrum damit um eine dritte DLT-Familie und um den Ressourcenverbrauch als Kriterium.

\paragraph{Architekturen ohne ausgewiesenes Identitätsmodell}
\label{par:literatur_ohne_idmodell}

Vier Arbeiten beschreiben keinen eigenständigen Mechanismus zur Identitätsverwaltung. Drei von ihnen sind vor allem wegen ihres Speicher- oder Evaluationsmodells von Interesse: Giao et al.~\cite{giao2026} verankern ausschließlich Metadaten -- Datentyp, Hash, Zeitstempel und Sensitivitätsflag -- auf einem privaten Ethereum-Netzwerk und belassen die Rohdaten vollständig off-chain, eine datenschutzrechtlich begründete Entscheidung, die direkt auf die Frage übertragbar ist, ob Identitätsobjekte in HCS-Nachrichten kodiert oder nur per Hash verankert werden. Gupta und Verma~\cite{guptaverma2026} erweitern die Sicherheitsdimension um privacy-erhaltende Zugriffskontrolle über Zero-Knowledge-Proofs und quantifizieren den Zielkonflikt zwischen Datenschutzniveau und Datennutzbarkeit. Amanat et al.~\cite{amanat2022} isolieren in ihrem Vergleich den Einfluss des Konsensmechanismus auf die Performanz, was für die Wahl der Vergleichsgrößen dieser Arbeit instruktiv ist.

Gesondert hervorzuheben ist Verma et al.~\cite{verma2024}: Die Arbeit schlägt die erste konkrete hashgraphbasierte Architektur für ein Managementsystem elektronischer Patientenakten vor und ist damit im Gesundheitskontext die einzige des Korpus, die nicht auf einer Blockchain aufsetzt. Verwendet wird ein Hedera-Netzwerk mit aBFT-Konsens; die Zugriffskontrolle erfolgt ticketbasiert über Smart Contracts, ohne DID- oder VC-Mechanismen. Die Autor:innen bezeichnen das Netzwerk als permissioned, betreiben ihre Smart Contracts jedoch auf dem öffentlich lesbaren Hedera-Netzwerk; in der Taxonomie wird die Arbeit daher dem öffentlichen Netzwerktyp zugeordnet. Die Evaluation vergleicht Transaktionszeiten mit einer Proof-of-Work-basierten Referenzimplementierung~\cite{verma2024}. Für diese Arbeit ist sie der zentrale Beleg dafür, dass Hashgraph im Gesundheitsdatenmanagement grundsätzlich tragfähig ist.

\paragraph{Zusammenfassung}
\label{par:literatur_architectures_zusammenfassung}

Über alle Arbeiten dieses Blocks hinweg verdichten sich drei wiederkehrende Architekturentscheidungen zu Vergleichsachsen für diese Arbeit. Erstens das Speichermuster: Die tatsächlichen Gesundheitsdaten liegen durchgängig off-chain, während der Ledger als Vertrauens- und Verifikationsanker Hashes, Metadaten, Zugriffsrichtlinien und Audit-Informationen verwaltet -- von Mamoru et al.~\cite{mamoru2026} mit Skalierbarkeitsaspekten und von Giao et al.~\cite{giao2026} mit dem Datenschutz begründet. Zweitens die Konsensauswahl als Kompromiss zwischen Fehlertoleranzmodell, Latenz, Finalität und Ressourcenaufwand, die von BFT-Varianten über Raft und Proof-of-Authority bis zu Proof-of-Stake reicht~\cite{joshimahajan2025, farin2025, amanat2022, tawfik2025}. Drittens die Verankerung von Zugriffskontrolle und Widerruf: Da der Ledger selbst nur Ordnungs- und Integritätsgarantien liefert, müssen Zugriffslogik, Einwilligungsmanagement und Widerruf auf Anwendungsebene implementiert werden~\cite{mamoru2026, tawfik2025, guptaverma2026}. Diese drei Achsen sind unmittelbar auf die in dieser Arbeit zu treffende Entscheidung übertragbar.


\subsection{Vergleichs- und Methodikliteratur}
\label{sec:literatur_vergleichsmethodik}

Der zweite Literaturblock umfasst Arbeiten, die in erster Linie methodische Beiträge zum systematischen Vergleich und zur Evaluation von DLT-Architekturen im Gesundheitswesen liefern und somit eine methodische Grundlage für den in dieser Arbeit vorgenommenen Vergleich von HCS und Hyperledger Indy in Bezug auf Performanz, Skalierbarkeit und Sicherheit. Da diese Arbeiten selbst keine eigene Architektur vorschlagen, werden sie nicht in die Taxonomie in Abschnitt~\ref{sec:literatur_architekturvergleich} aufgenommen.

Erler et al. entwickeln eine Taxonomie blockchainbasierter Anwendungen für das Gesundheitsdatenmanagement und leiten daraus ein Entscheidungsmodell für deren Entwurf ab \cite{erler2024}. Die Taxonomie folgt einem iterativen Entwicklungsverfahren von Nickerson et al. \cite{nickerson2013} und erweitert eine frühere Version um Dimensionen des Identitäts- und Zugriffsmanagements, darunter die Unterscheidung zwischen zentralisierten, föderierten, nutzerzentrierten und dezentralen Identitätsmanagementsystemen. Im Rahmen der dezentralen Kategorie unterscheiden Erler et al. zwischen Self-Sovereign Identity (SSI), bei der die Registrierung ohne Rückgriff auf eine bestehende Identität erfolgen kann, und Decentralized Trusted Identity (DTI), bei der die dezentrale Identität an eine bestehende Drittidentität gebunden wird \cite{erler2024}. Die Designentscheidungen des Modells beinhalten auch Annahmen über das Vertrauen der beteiligten Akteure, die mittels der STRIDE-Methode strukturiert analysiert werden. Für diese Arbeit ist die Taxonomie besonders relevant. Erstens liefert sie das methodische Vorgehen und die Klassifikationsprinzipien, an denen sich die entwickelte Taxonomie orientiert. Zweitens beschränkt sie sich nur auf blockchainbasierte Ansätze; hashgraphbasierte Architekturen sowie die zugrunde liegende Konsens-Familie werden nicht als Dimensionen erfasst. Genau diese technologische Achse ergänzt diese Arbeit.

Mazlan et al. \cite{mazlan2020} stellen eine systematische Literaturübersicht zu Skalierbarkeitsherausforderungen in blockchainbasierten Gesundheitssystemen vor. Als zentrale Herausforderungen heben die Autor:innen Blockgröße, hohes Datenvolumen, Transaktionsaufkommen, Knotenanzahl und Protokollbeschränkungen hervor. Die in der Literatur vorgeschlagenen 16 Lösungsansätze werden in zwei Hauptkategorien eingeordnet: Speicheroptimierung und Neugestaltung der Blockchain (u. a. Konsensprotokolle, Sharding, Lese- und Schreibmechanismen) \cite{mazlan2020}. Die Arbeit liefert damit eine Taxonomie von Skalierbarkeitsproblemen, aus der sich Vergleichskriterien für die Skalierbarkeit dieser Arbeit ableiten lassen.

Sehimi und Bendaoud \cite{sehimibendaoud2023} vertiefen diese Perspektive, indem sie konkrete Skalierbarkeitslösungen in blockchainbasierten EHR-Systemen betrachten, die in Off-Chain-Speicherung, Sidechains, Sharding und Konsensoptimierung unterteilt sind. Als Standardmetriken der Skalierbarkeitsevaluation schlagen die Autor:innen Durchsatz (TPS) und Latenz vor und veranschaulichen die Diskrepanz zwischen Blockchain-Systemen und traditionellen Netzwerken am Beispiel von Bitcoin (ca. 27 TPS) gegenüber dem Visa-Netzwerk (ca. 65.000 TPS) \cite{sehimibendaoud2023}. Die wichtigste Schlussfolgerung ist, dass viele vorgeschlagene EHR-Systeme ihre Skalierbarkeit überhaupt nicht oder nicht ausreichend evaluieren und sich hauptsächlich auf Off-Chain-Speicherung als alleinige Lösung verlassen, während Alternativen wie Sharding und Sidechains nicht ausreichend dargestellt bleiben. Die Autor:innen fordern daher einheitliche, durchsatz- und latenzbasierte Evaluationsstandards \cite{sehimibendaoud2023}, was eine Forderung ist, der die Evaluationsmethodik dieser Arbeit folgt.

Hashim et al. \cite{hashim2021} evaluieren die Performanz verschiedener Konsensalgorithmen (Proof-of-Work, PBFT sowie zwei Proof-of-Authority-Varianten - Aura, Apla) im Kontext des EHR-Austauschs und stellen zusätzlich ihren eigenen Echtzeit-PoA-Algorithmus vor. Die Simulation mit 100 bis 400 Knoten misst die Konsenszeit, Blockgenerierung, Durchsatz und Skalierungsverhalten und zeigt u. a., dass die Konsenszeit von PBFT linear mit der Knotenanzahl wächst, während PoW aufgrund des Rechenaufwands die höchste Konsenszeit aufweist \cite{hashim2021}. Für diese Arbeit ist diese Studie doppelt wichtig: sie liefert zum einen ein methodisches Modell für die simulations- bzw. messbasierte Evaluation von Konsensmechanismen, zum anderen konkrete Referenzwerte für die Einordnung der Konsenseigenschaften von Hedera Hashgraph (aBFT, Virtual Voting) im Vergleich zu in Indy-artigen permissioned Netzwerken verbreiteten BFT-Varianten.

Carlos Ferreira et al. \cite{ferreira2024} ergänzen den methodischen Rahmen mit einer PRISMA-basierten systematischen Literaturübersicht zu Interoperabilität und Sicherheit von EHR-Systemen auf DLT-Basis. Neben der Identifikation wichtigster Herausforderungen (fragmentierte Systeme, heterogene Standards, Datenschutzanforderungen) betonen die Autor:innen, dass die Interoperabilität zwischen verschiedenen Blockchain-Plattformen selbst eine offene Forschungsfrage ist und dass DLT-Lösungen mit bestehenden Standards und Protokolle des Gesundheitswesens kompatibel sein müssen \cite{ferreira2024}. Diese Perspektive ist für diese Arbeit relevant, da der Austausch der Ledger-Schicht in der von Erler et al. vorgeschlagenen Architektur \cite{erler2023} genau diesen plattformübergreifenden Ansatz erfordert.

Aus diesem Block ergibt sich für diese Arbeit eine methodische Grundlage: Durchsatz und Latenz als Kernmetriken der Performanz- und Skalierbarkeitsevaluation (\cite{sehimibendaoud2023, hashim2021, mazlan2020}), die Forderung nach expliziter, standardisierter Evaluation statt rein konzeptioneller Vorschläge (\cite{sehimibendaoud2023}) sowie Sicherheitsanforderungen als qualitative Vergleichsdimensionen (\cite{ferreira2024}).


\subsection{Ergänzende Literatur außerhalb des Gesundheitskontexts}
\label{sec:literatur_ergaenzend}

Außerhalb des Gesundheitskontexts gibt es erste Arbeiten, die Hashgraph-basierte DLT-Systeme als Grundlage für das dezentrale Identitätsmanagement untersuchen und damit die Übertragbarkeit des in dieser Arbeit verfolgten Ansatzes bestätigen.

Satybaldy et al. stellen die erste systematische Benchmark-Studie ledgerbasierter DID-Methoden vor und vergleichen \texttt{did:ethr} (Ethereum), \texttt{did:xrpl} (XRP Ledger) und \texttt{did:hedera} (Hedera) unter einer einheitlichen Installation anhand von Latenz, Transaktionskosten und On-Chain-Metadaten-Exposition \cite{satybaldy2026}. Die Autoren zeigen, dass keine Plattform alle drei Dimensionen gleichzeitig optimiert, bestätigen jedoch, dass Hedera die geringste On-Chain-Latenz bei geringster Metadaten-Exposition sowie die reifste Integration in bestehende SSI-Werkzeugketten hat, z.B. durch AnonCreds-Unterstützung über Aries-basierte Frameworks und eine vollständige Anbindung an den Universal Resolver \cite{satybaldy2026}. Für diese Arbeit ist die Studie besonders relevant, weil sie mit \texttt{did:hedera} exakt die DID-Methode evaluiert, auf der der umgesetzte Prototyp aufbaut, und eine methodische Vorlage für den Entwurf der Bewertungsdimensionen Performanz und Sicherheit auf Ebene einzelner DID-Lebenszyklusoperationen liefert. Da diese Arbeit weder den Gesundheitskontext noch die Architekturebene adressiert, wird sie als ergänzende Literatur eingeordnet.

Cádiz et al. demonstrieren die Anwendbarkeit hashgraphbasierter Konsensverfahren für dezentrales Identitätsmanagement im Kontext spezieller Transportnetze (VANETs) \cite{cadiz2025}. Die Autoren verwenden mit Teegraph einen hashgraphbasierten Konsensalgorithmus, der u.a. ein DAG-basiertes Gossip-Protokoll verwendet, und verankern W3C-konforme DIDs und VCs auf einem permissioned Teegraph-Ledger, um Fahrzeugauthentifizierung ohne wiederholte Anfragen an eine zentrale Trusted Authority zu ermöglichen \cite{cadiz2025}. Konzeptionell entspricht dies dem Schema, eine zentrale Vertrauensinstanz durch DLT-verankerte, dezentral verifizierbare Identitätsnachweise zu ersetzen. Die Arbeit belegt damit, dass die Kombination aus Hashgraph-Konsens und SSI-Prinzipien außerhalb des Gesundheitswesens tragfähig ist, bleibt jedoch auf einen Forschungsprototyp in einem anderen Anwendungsbereich beschränkt und wird daher ebenfalls als ergänzende Literatur klassifiziert.

\section{Taxonomische Einordnung der relevanten Paper}
\label{sec:literatur_architekturvergleich}

Um die im Abschnitt~\ref{sec:literatur_architectures} beschriebenen Architekturen systematisch zu vergleichen, werden die grundlegenden Architekturentscheidungen jeder Arbeit in einer Taxonomie (Tabelle~\ref{tab:taxonomie_core}) eingeordnet. Die Einordnung erfolgt entlang von sechs Kategorien, die verschiedenen Abstraktionsebenen zugeordnet werden können. Das Identitätsmodell bildet dabei die erste Spaltengruppe, da die Zeilenblöcke der Tabelle nach seinen Umsetzungen gegliedert sind und die Trennlinien zwischen den Blöcken damit der zugehörigen Spaltengruppe zugeordnet werden können. Der Anwendungskontext unterscheidet, ob eine Arbeit ihre Architektur im Gesundheitswesen oder in einer anderen Anwendungsdomäne verortet; er ist erforderlich, weil mit Cádiz et al. \cite{cadiz2025} eine Arbeit außerhalb des Gesundheitskontexts aufgenommen wird, um die im Gesundheitskontext fehlende Kombination aus dezentralem Identitätsmodell und hashgraphbasiertem Ledger im direkten Vergleich offensichtlich zu machen. Die Kategorie DLT-Art fasst die konkreten Frameworks zu Technologiefamilien zusammen und dient damit primär der Sichtbarmachung der Technologieverteilung im Literaturkorpus; für die plattformspezifische Einordnung bleibt die Kategorie Framework maßgeblich, die die konkret eingesetzte Plattform benennt. Der Netzwerktyp unterscheidet danach, wer die konsensbildenden Knoten betreibt: private Netzwerke unter der Kontrolle einer einzelnen Organisation, Konsortium-Netzwerke mehrerer gleichberechtigter Organisationen sowie öffentliche Netzwerke, deren Knotenmenge nicht von den teilnehmenden Organisationen bestimmt wird~\cite{ferdous2020}. Der Lesezugang ist von dieser Eigenschaft unabhängig und wird deshalb nicht als Ausprägung des Netzwerktyps geführt, sondern dort behandelt, wo er entscheidungsrelevant wird (Abschnitt~\ref{subsec:e1-netzwerkmodell}). Die Konsens-Familie ordnet den verwendeten Konsensmechanismus einer übergeordneten Familie zu. Das Identitätsmodell folgt der Unterscheidung von Erler et al. \cite{erler2024} zwischen dezentralen Modellen - Self-Sovereign Identity (SSI), bei der die Registrierung ohne Rückgriff auf eine bestehende Identität erfolgt, und Decentralized Trusted Identity (DTI), bei der die dezentrale Identität an eine bestehende Drittidentität gebunden wird - und schlüssel- bzw. zertifikatsbasierten Modellen, bei denen Identitäten von einer Betreiberinstanz ausgestellt und verwaltet werden. Das Zugriffsmodell sowie das On-/Off-Chain-Speichermodell der einzelnen Arbeiten wurden bereits in den Beschreibungen im Abschnitt~\ref{sec:literatur_architectures} dargestellt. Die Zeilen der Tabelle folgen derselben Gruppierung wie die Beschreibungen in Abschnitt~\ref{sec:literatur_architectures}: die vier durch Linien getrennten Blöcke entsprechen den Ausprägungen des Identitätsmodells, wobei der letzte Block die Arbeiten ohne ausgewiesenes Identitätsmodell zusammenfasst.

Die Dimensionen und Ausprägungen der Taxonomie wurden gemäß dem Taxonomie-Entwicklungsverfahren von Nickerson et al. \cite{nickerson2013} iterativ aus den in Abschnitt~\ref{sec:literatur_architectures} vorgestellten Architekturarbeiten abgeleitet, analog zum Vorgehen von Erler et al. \cite{erler2024}. Dabei wurde das Prinzip ihrer Klassifizierung übernommen: wo Publikationen einzelne Architekturaspekte nicht eindeutig beschreiben, erfolgt die Zuordnung allgemeiner oder wird ausgelassen, um annahmebasierte Klassifikationen zu vermeiden \cite{erler2024}. Dies betrifft vor allem die Konsens-Familie und den Netzwerktyp: falls eine Arbeit den eingesetzten Konsensmechanismus oder den Netzwerktyp nicht explizit erwähnt, bleibt die entsprechende Zeile unbesetzt, auch wenn die verwendete Plattform eine Voreinstellung vorsieht.


\begin{landscape}
\begin{table}[htbp]
\centering
\scriptsize
\renewcommand{\arraystretch}{1.35}
\setlength{\tabcolsep}{3pt}
\caption{Taxonomie der relevanten Paper: architektonische Einordnung}
\label{tab:taxonomie_core}
\begin{tabular}{p{4.2cm} ccc cc ccc ccc cccccccc cccccc}
\toprule
& \multicolumn{3}{c}{\textbf{Identitätsmodell}} & \multicolumn{2}{c}{\textbf{Anw.-kontext}} & \multicolumn{3}{c}{\textbf{DLT-Art}} & \multicolumn{3}{c}{\textbf{Netzwerktyp}} & \multicolumn{8}{c}{\textbf{Framework}} & \multicolumn{6}{c}{\textbf{Konsens-Familie}} \\
\cmidrule(lr){2-4} \cmidrule(lr){5-6} \cmidrule(lr){7-9} \cmidrule(lr){10-12} \cmidrule(lr){13-20} \cmidrule(lr){21-26}
\textbf{Paper} &
\rotatebox{90}{Dezentral: SSI} & \rotatebox{90}{Dezentral: DTI} & \rotatebox{90}{Schlüssel-/zertifikatsb.} &
\rotatebox{90}{Gesundheitswesen} & \rotatebox{90}{Andere Domäne} &
\rotatebox{90}{Blockchain} & \rotatebox{90}{Hashgraph} & \rotatebox{90}{Corda} &
\rotatebox{90}{Privat} & \rotatebox{90}{Konsortium} & \rotatebox{90}{Öffentlich} &
\rotatebox{90}{Hyperledger Indy/Aries} & \rotatebox{90}{Hyperledger Fabric} & \rotatebox{90}{Hyperledger Besu} & \rotatebox{90}{Ethereum (privat)} & \rotatebox{90}{Hedera} & \rotatebox{90}{R3 Corda} & \rotatebox{90}{EBSI} & \rotatebox{90}{Teegraph} &
\rotatebox{90}{BFT (Voting)} & \rotatebox{90}{aBFT (Gossip)} & \rotatebox{90}{Leader-b. (Raft)} & \rotatebox{90}{Proof of Authority} & \rotatebox{90}{Proof of Stake} & \rotatebox{90}{Notary-basiert} \\
\midrule
Erler et al. \cite{erler2023} & \checkmark &    &    & \checkmark &    & \checkmark &    &    &    &    & \checkmark & \checkmark &    &    &    &    &    &    &    & \checkmark &    &    &    &    &    \\
Manoj et al. \cite{manoj2022} & \checkmark &    &    & \checkmark &    & \checkmark &    &    &    & \checkmark &    & \checkmark &    &    &    &    &    &    &    & \checkmark &    &    &    &    &    \\
Torongo, Toorani \cite{torongo2023} & \checkmark &    &    & \checkmark &    & \checkmark &    &    &    &    & \checkmark & \checkmark &    &    &    &    &    &    &    & \checkmark &    &    &    &    &    \\
Chainarong et al. \cite{chainarong2025} & \checkmark &    &    & \checkmark &    & \checkmark &    &    &    & \checkmark &    &    & \checkmark &    &    &    &    &    &    &    &    &    &    &    &    \\
Cádiz et al. \cite{cadiz2025} & \checkmark &    &    &    & \checkmark &    & \checkmark &    &    &    &    &    &    &    &    &    &    &    & \checkmark &    & \checkmark &    &    &    &    \\
\midrule
Bittins et al. \cite{bittins2021} &    & \checkmark &    & \checkmark &    & \checkmark &    &    &    & \checkmark &    &    & \checkmark &    &    &    &    & \checkmark &    &    &    &    &    &    &    \\
\midrule
Tawfik et al. \cite{tawfik2025} &    &    & \checkmark & \checkmark &    & \checkmark &    &    &    & \checkmark &    &    & \checkmark &    &    &    &    &    &    & \checkmark &    & \checkmark &    &    &    \\
Joshi, Mahajan \cite{joshimahajan2025} &    &    & \checkmark & \checkmark &    & \checkmark &    &    &    & \checkmark &    &    & \checkmark &    &    &    &    &    &    &    &    & \checkmark &    &    &    \\
Mamoru et al. \cite{mamoru2026} &    &    & \checkmark & \checkmark &    & \checkmark &    &    &    & \checkmark &    &    & \checkmark &    &    &    &    &    &    &    &    &    &    &    &    \\
Farin et al. \cite{farin2025} &    &    & \checkmark & \checkmark &    & \checkmark &    &    & \checkmark &    &    &    &    & \checkmark &    &    &    &    &    &    &    &    & \checkmark &    &    \\
Minango et al. \cite{minango2025} &    &    & \checkmark & \checkmark &    &    &    & \checkmark & \checkmark &    &    &    &    &    &    &    & \checkmark &    &    &    &    &    &    &    & \checkmark \\
\midrule
Giao et al. \cite{giao2026} &    &    &    & \checkmark &    & \checkmark &    &    & \checkmark &    &    &    &    &    & \checkmark &    &    &    &    &    &    &    & \checkmark &    &    \\
Gupta, Verma \cite{guptaverma2026} &    &    &    & \checkmark &    & \checkmark &    &    &    & \checkmark &    &    & \checkmark &    &    &    &    &    &    &    &    &    &    &    &    \\
Amanat et al. \cite{amanat2022} &    &    &    & \checkmark &    & \checkmark &    &    &    &    &    &    &    &    &    &    &    &    &    &    &    &    &    & \checkmark &    \\
Verma et al. \cite{verma2024} &    &    &    & \checkmark &    &    & \checkmark &    &    &    & \checkmark &    &    &    &    & \checkmark &    &    &    &    & \checkmark &    &    &    &    \\
\bottomrule
\end{tabular}
\end{table}
\end{landscape}


Die Taxonomie macht zwei Gliederungen offensichtlich. Erstens die Technologieverteilung: zwölf der fünfzehn in der Taxonomie erfassten Arbeiten setzen auf eine Blockchain, davon drei auf Hyperledger Indy mit einer Aries-Agentenschicht und sechs auf Hyperledger Fabric; zwei Arbeiten verwenden einen hashgraphbasierten \cite{verma2024, cadiz2025} und eine einen Corda-basierten Ledger \cite{minango2025}. Zweitens die Verteilung der Identitätsmodelle: Fünf Arbeiten weisen ein SSI-basiertes und eine Arbeit ein DTI-basiertes Modell aus, fünf Arbeiten verwalten Identitäten schlüssel- bzw. zertifikatsbasiert, und vier Arbeiten beschreiben kein eigenständiges Identitätsmodell. Erst die Überlagerung beider Verteilungen ist für diese Arbeit aussagekräftig; sie ist ein Gegenstand der folgenden Schlussfolgerungen.

\section{Schlussfolgerungen: Forschungslücke und Vergleichskriterien}
\label{sec:literatur_schlussfolgerungen}

Aus der Taxonomie im Abschnitt~\ref{sec:literatur_architekturvergleich} können drei grundlegende Schlussfolgerungen abgeleitet werden, auf denen die Forschungslücke sowie die Vergleichskriterien dieser Arbeit basieren.

\paragraph{Schlussfolgerung 1: Konzentration auf blockchainbasierte Ledger.}
Die Tabelle~\ref{tab:taxonomie_core} zeigt eine klare Konzentration der Literatur auf blockchainbasierte Ledger-Technologien, insbesondere aus dem Hyperledger-Ökosystem: Zwölf der vierzehn im Gesundheitskontext untersuchten Arbeiten basieren auf einer Blockchain, davon neun auf Hyperledger Indy oder Fabric. Nur zwei dieser Arbeiten weichen von diesem Muster ab: Verma et al. \cite{verma2024} mit einem hashgraphbasierten und Minango et al. \cite{minango2025} mit einem Corda-basierten Ansatz. Die einzige weitere hashgraphbasierte Arbeit der Taxonomie, Cádiz et al. \cite{cadiz2025}, ist in der Spaltengruppe Anwendungskontext einer anderen Domäne zugeordnet und zählt damit nicht zu diesem Korpus. Alternative DLT-Familien sind im Gesundheitskontext daher deutlich unterrepräsentiert, obwohl ihnen in der technologieübergreifenden Vergleichsliteratur (Abschnitt~\ref{sec:literatur_vergleichsmethodik}) günstige Durchsatz-, Latenz- und Fairness-Eigenschaften zugeschrieben werden.

\paragraph{Schlussfolgerung 2: Keine hashgraphbasierte SSI-Architektur im Gesundheitskontext.}
Die Spaltengruppen Identitätsmodell, DLT-Art und Anwendungskontext der Tabelle~\ref{tab:taxonomie_core} machen zwei bisher unverbundene Entwicklungslinien sichtbar. Alle Arbeiten des Gesundheitskontexts mit einem dezentralen Identitätsmodell - sowohl die SSI-basierten Ansätze von Erler et al. \cite{erler2023}, Manoj et al. \cite{manoj2022}, Torongo und Toorani \cite{torongo2023} sowie Chainarong et al. \cite{chainarong2025} als auch der DTI-basierte Ansatz von Bittins et al. \cite{bittins2021} - basieren auf blockchainbasierten Ledgern. Umgekehrt vermeidet die einzige hashgraphbasierte Architektur des Gesundheitskontexts von Verma et al. \cite{verma2024} vollständig ein dezentrales Identitätsmodell und verwendet stattdessen einen ticketbasierten Zugriffsmechanismus. Die einzige Zeile der Taxonomie, die ein SSI-basiertes Identitätsmodell mit einem hashgraphbasierten Ledger verbindet, ist die von Cádiz et al. \cite{cadiz2025} - und genau diese Zeile ist in der Spaltengruppe Anwendungskontext einer anderen Domäne zugeordnet. Die Überlagerung der drei Spaltengruppen macht damit direkt sichtbar, dass eine Architektur, die beide Entwicklungslinien im Gesundheitskontext zusammenführt, in der betrachteten Literatur nicht existiert, obwohl die Kombination als solche technisch tragfähig ist \cite{cadiz2025, satybaldy2026}.

\paragraph{Schlussfolgerung 3: Eingeschränkte Vergleichbarkeit der publizierten Evaluationsergebnisse.} 
Die Evaluationsansätze der untersuchten Arbeiten unterscheiden sich sehr: Sie reichen von rein qualitativen Anforderungsanalysen \cite{erler2023, bittins2021} über Latenz- und Durchsatzmessungen einzelner Workflows \cite{manoj2022, torongo2023, chainarong2025} bis hin zu systematischen Benchmarks mit dedizierten Werkzeugen und Angriffssimulationen \cite{tawfik2025, joshimahajan2025}. Die berichteten Kennzahlen beziehen sich damit auf unterschiedliche Messgrößen, Lastprofile und Systemgrenzen und sind untereinander nicht vergleichbar; die methodische Literatur bestätigt dies insbesondere für Skalierbarkeitseigenschaften, die häufig gar nicht oder inkonsistent evaluiert werden (Abschnitt~\ref{sec:literatur_vergleichsmethodik}). Ein Vergleich der Ledger-Technologien lässt sich daher nicht aus der vorhandenen Literatur ableiten, sondern erfordert eine eigene Messung beider Varianten unter gleichen Bedingungen.

\paragraph{Forschungslücke.}
Die drei Schlussfolgerungen ergeben eine Forschungslücke, die diese Arbeit adressiert: es fehlt sowohl eine Architektur, die ein SSI/DID-basiertes Identitätsmodell im Gesundheitskontext auf einem hashgraphbasierten Ledger realisiert, als auch ein systematischer, empirisch fundierter Vergleich zwischen einem hashgraphbasierten und einem blockchainbasierten Ledger bei einer ansonsten unveränderten Anwendungsarchitektur. Die in Abschnitt~\ref{sec:literatur_ergaenzend} beschriebenen Arbeiten \cite{satybaldy2026, cadiz2025} zeigen, dass die Kombination aus Hashgraph-Konsens und dezentralem Identitätsmodell technisch tragfähig ist; sie adressieren jedoch weder den Gesundheitskontext noch die Architekturebene und schließen die Lücke damit nicht. Genau diese Lücke füllt diese Arbeit, indem die Ledger-Schicht der Referenzarchitektur von Erler et al. \cite{erler2023} durch eine auf dem Hedera Consensus Service basierende Alternative ersetzt und beide Varianten entlang einheitlicher Kriterien evaluiert werden. Das Ersetzen nur der Ledger-Schicht bei unveränderter Agenten- und Kommunikationsschicht ermöglicht dabei eine isolierte Betrachtung des Technologieeinflusses.

\paragraph{Abgeleitete Vergleichskriterien.}
Auf der Grundlage der Evaluationsansätze der betrachteten Architekturen (Abschnitt~\ref{sec:literatur_architectures}) und der Metriken der methodischen Literatur (Abschnitt~\ref{sec:literatur_vergleichsmethodik}) werden die Vergleichskriterien für diese Arbeit durch drei Dimensionen Performanz, Skalierbarkeit und Sicherheit abgeleitet. Die Tabelle~\ref{tab:synthese_vergleichskriterien} fasst die Kriterien, ihre Beschreibung und ihre Herleitung aus der Literatur zusammen.

\begin{table}[htbp]
\centering
\scriptsize
\renewcommand{\arraystretch}{1.15}
\caption{Synthese der Vergleichskriterien entlang der Dimensionen
Performanz, Skalierbarkeit und Sicherheit}
\begin{tabular}{p{2.2cm}p{3.6cm}p{4.6cm}p{2.8cm}}
\toprule
\textbf{Dimension} & \textbf{Kriterium} & \textbf{Beschreibung} &
\textbf{Herleitung} \\
\midrule
Performanz & Durchsatz & Transaktionen pro Sekunde unter definierter Last &
\cite{tawfik2025, verma2024, minango2025, hashim2021} \\
& Latenz & Zeit von Transaktionseinreichung bis Finalität, differenziert
nach Operationstyp (z.\,B. DID-Registrierung, Credential-Prüfung,
Widerruf) & \cite{manoj2022, torongo2023, mamoru2026, hashim2021} \\
& Ressourcenverbrauch & CPU- und Speicherauslastung der beteiligten
Komponenten & \cite{minango2025, amanat2022} \\
\midrule
Skalierbarkeit & Lastskalierung & Entwicklung von Durchsatz und Latenz bei
steigender Transaktionsrate und Nutzerzahl & \cite{joshimahajan2025,
chainarong2025, mazlan2020, sehimibendaoud2023} \\
& Datenwachstum & Entwicklung des Ledger-Speicherbedarfs bei wachsendem
Datenbestand (On-Chain-Fußabdruck) & \cite{giao2026, mamoru2026} \\
\midrule
Sicherheit & Konsens-Fehlertoleranz & Toleranz gegenüber byzantinischen
Knoten und Finalitätsgarantien des Konsensmechanismus (aBFT vs.
RBFT/PBFT/Raft/PoA) & \cite{verma2024, tawfik2025, farin2025} \\
& Zugriffskontrolle und Widerruf & Durchsetzbarkeit von
Zugriffsrichtlinien und Wirksamkeit des Credential-Widerrufs &
\cite{torongo2023, mamoru2026} \\
\bottomrule
\end{tabular}
\label{tab:synthese_vergleichskriterien}
\end{table}
\normalsize

\section{Referenzarchitektur nach Erler et al.}
\label{sec:sota_referenzarchitektur}
Die im Abschnitt~\ref{sec:literatur_architectures} eingeordnete Arbeit von Erler et al. \cite{erler2023} spielt für diese Arbeit eine besondere Rolle: Sie ist nicht nur eine von vierzehn verwandten Arbeiten im Gesundheitskontext, sondern der konkrete Gegenstand, dessen Ledger-Schicht im Rahmen des Proof-of-Concepts ersetzt wird. Daher wird sie an dieser Stelle im Detail dargestellt. Die von Erler et al. vorgeschlagene Architektur einer dezentralen Anwendung zur Verwaltung von Gesundheitsdaten (Health Data Management Application, HDMA) wurde systematisch auf der Grundlage einer Bedrohungsanalyse entworfen: die Autoren identifizierten schützenswerte Assets und Vertrauensgrenzen basierend auf dem Flussdiagramm des Zielsystems, analysierten potenzielle Bedrohungen mit der STRIDE-Methode und leiteten daraus Schutzmaßnahmen ab, die in Expertenworkshops in Anforderungen zusammengefasst und schließlich in eine Systemarchitektur überführt wurden \cite{erler2023}. Daher ist die Architektur explizit sicherheitsgetrieben entworfen.

\subsection{Rollen und Komponenten}
Das System umfasst vier Nutzerrollen \cite{erler2023}: Patienten als Datensubjekte, die über eine mobile App auf ihre Gesundheitsdaten zugreifen und granuläre Zugriffsrichtlinien für einzelne Dokumente definieren; medizinische Fachkräfte, die medizinische Dokumente in den internen Krankenhausinformationssystemen (KIS) ihrer Einrichtung erstellen und - mit Einwilligung der Patienten - auch Daten anderer Einrichtungen verwenden; Angehörige, denen Patienten Lesezugriff oder Verwaltungsrechte erteilen können; sowie Systemadministratoren der medizinischen Einrichtungen, die für die Funktionsweise der Komponenten verantwortlich sind.

Die Komponentenebene des Systems besteht aus einer mobilen App der Patienten und Angehörigen, dem Netzwerk-Konnektor für medizinische Einrichtungen sowie der Verbindungs-Netzwerkinfrastruktur \cite{erler2023}. Der Netzwerk-Konnektor dient als Gateway zwischen dem internen KIS einer Einrichtung und dem Netzwerk: Er umfasst ein API-Gateway als zentrale Schnittstelle, eine Metadaten-Datenbank (u. a. für die Zuordnung von Peer-DIDs zu internen Nutzerkennungen, Zugriffsrichtlinien und Zugriffsprotokolle) sowie ein Web-Frontend für medizinischen Fachkräfte. Die von Patienten definierten Zugriffsrichtlinien werden auf dem Netzwerk-Konnektor angewendet. Ein zentrales Entwurfsprinzip ist dabei die Datenminimierung: die Gesundheitsdaten selbst bleiben in den internen Systemen der erzeugenden Einrichtung bzw. lokal auf den mobilen Endgeräten (off-chain) - es entstehen damit keine redundanten zentralen Datensilos \cite{erler2023}.

\subsection{Netzwerkinfrastruktur: DIDComm-Netzwerk aus Aries-Agenten}
Die Netzwerkinfrastruktur ist auf Basis des DIDComm-Protokolls als dezentrales Peer-to-Peer-Kommunikationsnetz konzipiert und basiert auf Open-Source-Implementierungen Hyperledger Aries und Indy \cite{erler2023}. Jeder Teilnehmer wird durch einen Aries-Agenten vertreten (siehe Abschnitt~\ref{sec:grundlagen-blockchain}): Die mobilen Apps integrieren Aries Mobile Agents, die Netzwerk-Konnektoren der Einrichtungen verwenden Aries Cloud Agents; zusammen bilden sie die Edge-Agenten des Netzwerks. Da mobile Endgeräte nicht dauerhaft verfügbar sind, funktionieren Aries Mediator Agents als eine Art Briefkasten, das eine asynchrone verschlüsselte Zustellung von Nachrichten an mobile Agenten ermöglicht \cite{erler2023}. Für jede private Verbindung zwischen den beiden Agenten wird ein paarweises Peer-DID erzeugt, das nur den zwei beteiligten Parteien bekannt ist, wodurch die Korrelierbarkeit von Beziehungen minimiert wird \cite{erler2023}.

\subsection{Rolle der Ledger-Schicht}
\label{subsec:sota_ledgerschicht}
Das DIDComm-Netzwerk wird durch eine public permissioned Indy-Blockchain als zusätzlichen Vertrauensanker ergänzt~\cite{erler2023}. Sie führt ausschließlich pseudonyme Identitätsobjekte -- öffentliche DIDs, Schemas, Credential Definitions und Widerrufsinformationen --, während die medizinischen Dokumente off-chain verbleiben.

Diese Zuweisung ist für die vorliegende Arbeit der entscheidende Anknüpfungspunkt, denn sie bestimmt, welche Funktionen eine alternative Implementierung der Ledger-Schicht erbringen muss, während die Agentenschicht unverändert bleibt. Die systematische Ableitung dieser Funktionen und ihrer Indy-spezifischen Umsetzung erfolgt in Abschnitt~\ref{subsec:ledger-funktionen} und wird hier nicht vorweggenommen.